\documentclass[compress, xcolor={svgnames}, aspectratio=169, smaller, t]{beamer}
\newcommand{\CommonPath}{../../Common}
\newcommand{\coursename}{Intro to DL - S9}
\usepackage{\CommonPath/slidespreamble}
\usepackage{\CommonPath/math_commands}
\setbeameroption{hide notes}

\title{Introduction to Deep Learning}
\subtitle{Session 9: Convolutional Neural Networks \\[1em]
\small May 2026}
\author{Magda Gregorová - \href{mailto:magda.gregorova@thws.de}{magda.gregorova@thws.de}}
\institute{
\includegraphics[width=0.3\textwidth]{thws-logo_vert_en_orange-rgb.png}}
\hypersetup{
pdftitle={Introduction to Deep Learning},
pdfauthor={Magda Gregorová},
pdfkeywords={convolutional neural networks, convolution, pooling, ResNet, AlexNet, VGG}
}
\date{{\color{skyblue1}\footnotesize Licensed according to
\href{https://creativecommons.org/licenses/by-sa/4.0}{\color{skyblue1}CC-BY-SA 4.0}}}

\begin{document}

\begin{frame}[plain]
    \titlepage
\end{frame}

\begin{frame}{Lecture Overview}
\tableofcontents
\end{frame}

% -------------------------------------------------------
\section{Why Not MLPs for Images?}
% -------------------------------------------------------

\begin{frame}{Images in Computers}

\vskip 1em

\begin{columns}[c]
\begin{column}{0.7\textwidth}

\structure{Image:} 3D tensor of shape $C \times H \times W$
\begin{itemize}
    \item Each entry: pixel intensity, integer 0--255 or float 0--1,
    \item $2^8$ possible values $\sim 8$ bits
    \item \structure{Grayscale} ($C=1$): single channel, 0 = black, 255 = white
    \item \structure{RGB} ($C=3$): red, green, blue --- standard cameras, screens
\end{itemize}

\end{column}
\begin{column}{0.25\textwidth}
\centering
\includegraphics[width=\textwidth]{\CommonPath/Pics/rgb.png}
\end{column}
\end{columns}

\structure{Other formats:}
\begin{itemize}
    \item RGBA ($C=4$): RGB + $\alpha$ transparency
    \item CMYK ($C=4$): print format
    \item Medical/hyperspectral: many channels, often 16-bit depth
\end{itemize}

\vskip 1em

\centering
\alert{DL mostly \textbf{RGB and grayscale} --- other formats typically converted before use}

\end{frame}

\begin{frame}{Images as Data}

\vskip 1em

\structure{Image sizes --- from benchmarks to real world:}

\begin{center}
\begin{tabular}{llll}
\toprule
\textbf{Source} & \textbf{Size} & \textbf{Channels} & \textbf{Pixels} \\
\midrule
MNIST & $28 \times 28$ & 1 & 784 \\
ImageNet & $224 \times 224$ & 3 & 150\,528 \\
Phone camera & $4000 \times 3000$ & 3 & 36\,000\,000 \\
Full HD screen & $1920 \times 1080$ & 3 & 6\,220\,800 \\
Medical (CT slice) & $512 \times 512$ & 1 & 262\,144 \\
\bottomrule
\end{tabular}
\end{center}

\vskip 1em

\structure{MLP approach:} flatten image to vector $\vx \in \mathbb{R}^{H \cdot W \cdot C}$, apply fully connected layers

\vskip 0.5em

ImageNet: first hidden from $\R^{150\,528} \to \R^{1024} \Rightarrow \approx 154$M parameters --- one layer only!

\vskip 2em

\centering
\alert{\large How to handle high-dimensional data efficiently?}

\end{frame}

\begin{frame}{How to Look at Images}

\vskip 2em

\begin{columns}[t]
\begin{column}{0.48\textwidth}
\structure{Human vision:}
\begin{itemize}
    \item Looks for local patterns --- edges, corners, textures
    \item Same pattern recognised anywhere in the visual field
    \item Simple features combine into complex ones
\end{itemize}
\end{column}
\begin{column}{0.48\textwidth}
\structure{Technical equivalent:}
\begin{itemize}
    \item \structure{Local connectivity} --- connect to small patches only
    \item \structure{Parameter sharing} --- same filter at every location
    \item \structure{Hierarchical features} --- edges $\to$ shapes $\to$ objects
\end{itemize}
\end{column}
\end{columns}

\vskip 2em

\begin{center}
Three key properties: \textbf{\alert{locality}, \alert{translation invariance}, \alert{parameter efficiency}}\\[1em]
These define the \textbf{convolutional layer}
\end{center}

\end{frame}

% -------------------------------------------------------
\section{Convolution Operation}
% -------------------------------------------------------

\begin{frame}{}
\vskip 8em
\begin{center}
{\Large \structure{Convolution Operation}}
\end{center}
\end{frame}

\begin{frame}{Convolution --- The Operation}

\vskip 1em

\centering

\structure{Single step:} dot product of local patch and filter $\to$ one output value\\[0.5em]
\tikzsetnextfilename{conv2d}
% conv2d.tex — placeholder
\begin{tikzpicture}[scale=0.7]
% input image
\foreach \i in {0,...,4} {
    \foreach \j in {0,...,4} {
        \draw[thwspetrol!40, fill=thwspetrol!10] (\j*0.5,\i*0.5) rectangle (\j*0.5+0.45,\i*0.5+0.45);
    }
}
% highlighted patch
\foreach \i in {1,...,3} {
    \foreach \j in {1,...,3} {
        \draw[thwsorange, fill=thwsorange!30, thick] (\j*0.5,\i*0.5) rectangle (\j*0.5+0.45,\i*0.5+0.45);
    }
}
\node[font=\small] at (1.1,-0.3) {input};
% filter
\foreach \i in {0,...,2} {
    \foreach \j in {0,...,2} {
        \draw[thwsorange, fill=thwsorange!50, thick] (\j*0.5+3.5,\i*0.5+0.5) rectangle (\j*0.5+3.95,\i*0.5+0.95);
    }
}
\node[font=\small] at (4.25,0.2) {filter};
% output
\foreach \i in {0,...,2} {
    \foreach \j in {0,...,2} {
        \draw[carnelian!60, fill=carnelian!15] (\j*0.5+6.5,\i*0.5+0.5) rectangle (\j*0.5+6.95,\i*0.5+0.95);
    }
}
\draw[carnelian, fill=carnelian!40, thick] (7.0,1.0) rectangle (7.45,1.45);
\node[font=\small] at (7.25,0.2) {output};
\draw[->, thick, thwsgrey] (3.0,1.25) -- (3.4,1.25);
\draw[->, thick, thwsgrey] (5.5,1.25) -- (6.4,1.25);
\end{tikzpicture}


\vskip 0.5em

\structure{Full operation:} same filter slides across input --- each position produces one output value\\[0.5em]
\tikzsetnextfilename{conv2d_steps}
% conv2d_steps.tex — three filter positions on same input
\begin{tikzpicture}[scale=0.72]

% --- position 1: top-left (patch at col 0-2, row 2-4 in bottom-up coords) ---
\begin{scope}[xshift=0cm]
    % full grid
    \foreach \col in {0,...,4} {
        \foreach \row in {0,...,4} {
            \draw[thwspetrol!25, fill=thwspetrol!7]
                (\col*0.52, \row*0.52) rectangle (\col*0.52+0.48, \row*0.52+0.48);
        }
    }
    % patch cols 0-2, rows 2-4 (bottom-up: rows 0-2)
    \foreach \col in {0,1,2} {
        \foreach \row in {2,3,4} {
            \draw[thwsorange, fill=thwsorange!30, thick]
                (\col*0.52, \row*0.52) rectangle (\col*0.52+0.48, \row*0.52+0.48);
        }
    }
    % output grid 3x3
    \foreach \col in {0,1,2} {
        \foreach \row in {0,1,2} {
            \draw[carnelian!30, fill=carnelian!8]
                (\col*0.52+2.9, \row*0.52+0.52) rectangle
                (\col*0.52+3.38, \row*0.52+1.0);
        }
    }
    % highlight top-left output
    \draw[carnelian, fill=carnelian!40, thick]
        (2.9, 1.54) rectangle (3.38, 2.02);
    \node[font=\scriptsize, thwsorange] at (2.7,-0.3) {step 1};

\draw[->, thick, thwsgrey] (4.7,1.3) -- (5.5,1.3);

\end{scope}

% --- position 2: shifted right (patch at col 1-3, row 2-4) ---
\begin{scope}[xshift=6cm]
    \foreach \col in {0,...,4} {
        \foreach \row in {0,...,4} {
            \draw[thwspetrol!25, fill=thwspetrol!7]
                (\col*0.52, \row*0.52) rectangle (\col*0.52+0.48, \row*0.52+0.48);
        }
    }
    \foreach \col in {1,2,3} {
        \foreach \row in {2,3,4} {
            \draw[thwsorange, fill=thwsorange!30, thick]
                (\col*0.52, \row*0.52) rectangle (\col*0.52+0.48, \row*0.52+0.48);
        }
    }
    \foreach \col in {0,1,2} {
        \foreach \row in {0,1,2} {
            \draw[carnelian!30, fill=carnelian!8]
                (\col*0.52+2.9, \row*0.52+0.52) rectangle
                (\col*0.52+3.38, \row*0.52+1.0);
        }
    }
    % highlight second output in top row
    \draw[carnelian, fill=carnelian!40, thick]
        (3.42, 1.54) rectangle (3.9, 2.02);
    \node[font=\scriptsize, thwsorange] at (2.7,-0.3) {step 2};

\draw[->, thick, thwsgrey] (4.7,1.3) -- (5.5,1.3);
\end{scope}

% --- position 3: shifted down (patch at col 0-2, row 1-3) ---
\begin{scope}[xshift=12cm]
    \foreach \col in {0,...,4} {
        \foreach \row in {0,...,4} {
            \draw[thwspetrol!25, fill=thwspetrol!7]
                (\col*0.52, \row*0.52) rectangle (\col*0.52+0.48, \row*0.52+0.48);
        }
    }
    \foreach \col in {0,1,2} {
        \foreach \row in {1,2,3} {
            \draw[thwsorange, fill=thwsorange!30, thick]
                (\col*0.52, \row*0.52) rectangle (\col*0.52+0.48, \row*0.52+0.48);
        }
    }
    \foreach \col in {0,1,2} {
        \foreach \row in {0,1,2} {
            \draw[carnelian!30, fill=carnelian!8]
                (\col*0.52+2.9, \row*0.52+0.52) rectangle
                (\col*0.52+3.38, \row*0.52+1.0);
        }
    }
    % highlight second output in first column
    \draw[carnelian, fill=carnelian!40, thick]
        (2.9, 1.04) rectangle (3.38, 1.52);
    \node[font=\scriptsize, thwsorange] at (2.7,-0.3) {step 3};
\end{scope}

% \node[font=\scriptsize, thwspetrol, align=center] at (8.5,-1)
%     {same filter slides across input --- each position produces one output value};

\end{tikzpicture}


\vskip 0.5em

\structure{All 3 requirements:} local connectivity, parameter sharing, translation invariance

\end{frame}

\begin{frame}{1D Convolution}

\vskip 2em

\structure{Filter $\kernel \in \mathbb{R}^k$} slides over signal $\vx \in \mathbb{R}^n$ --- at each position $i$, dot product with local patch

\begin{columns}[c]
\begin{column}{0.4\textwidth}
\centering
\tikzsetnextfilename{conv1d}
% conv1d.tex — placeholder
\begin{tikzpicture}[scale=0.8]
\foreach \i in {0,...,7} {
    \draw[thwspetrol, fill=thwspetrol!20] (\i*0.6,0) rectangle (\i*0.6+0.5,0.5);
    \node[font=\tiny] at (\i*0.6+0.25,0.25) {$x_{\i}$};
}
\foreach \i in {0,...,2} {
    \draw[thwsorange, fill=thwsorange!30, thick] (\i*0.6+0.6,0.8) rectangle (\i*0.6+1.1,1.3);
    \node[font=\tiny] at (\i*0.6+0.85,1.05) {$w_{\i}$};
}
\foreach \i in {0,...,5} {
    \draw[carnelian, fill=carnelian!20] (\i*0.6+0.6,-0.8) rectangle (\i*0.6+1.1,-0.3);
    \node[font=\tiny] at (\i*0.6+0.85,-0.55) {$y_{\i}$};
}
\draw[->, thwsgrey] (1.45,0.8) -- (1.45,-0.3);
\node[font=\tiny, thwsgrey] at (1.8,0.25) {$\cdot$};
\end{tikzpicture}

\end{column}
\begin{column}{0.55\textwidth}

\[
    (\vx * \kernel)[i] = \sum_{j=0}^{k-1} \vx[i+j]\; \kernel[j]
\]

\begin{itemize}
    \item $i$ --- output position; $j$ --- index within filter
    \item Dot product of filter with local patch
\end{itemize}

\vskip 1em

\structure{Example:} at position $i=1$:
\[
    1{\cdot}1 + 3{\cdot}(-1) + 4{\cdot}2 = 6
\]

\end{column}
\end{columns}

\end{frame}

\begin{frame}{2D Convolution}

\vskip 2em

\structure{Extend to 2D:} filter $\kernel \in \mathbb{R}^{k \times k}$ slides over image $\mX \in \mathbb{R}^{H \times W}$:
\[
    (\mX * \kernel)[i,j] = \sum_{p=0}^{k-1}\sum_{q=0}^{k-1} \mX[i+p,\, j+q]\; \kernel[p,q]
\]

\begin{columns}[c]
\begin{column}{0.5\textwidth}
\centering
\tikzsetnextfilename{conv2d}
% conv2d.tex — placeholder
\begin{tikzpicture}[scale=0.7]
% input image
\foreach \i in {0,...,4} {
    \foreach \j in {0,...,4} {
        \draw[thwspetrol!40, fill=thwspetrol!10] (\j*0.5,\i*0.5) rectangle (\j*0.5+0.45,\i*0.5+0.45);
    }
}
% highlighted patch
\foreach \i in {1,...,3} {
    \foreach \j in {1,...,3} {
        \draw[thwsorange, fill=thwsorange!30, thick] (\j*0.5,\i*0.5) rectangle (\j*0.5+0.45,\i*0.5+0.45);
    }
}
\node[font=\small] at (1.1,-0.3) {input};
% filter
\foreach \i in {0,...,2} {
    \foreach \j in {0,...,2} {
        \draw[thwsorange, fill=thwsorange!50, thick] (\j*0.5+3.5,\i*0.5+0.5) rectangle (\j*0.5+3.95,\i*0.5+0.95);
    }
}
\node[font=\small] at (4.25,0.2) {filter};
% output
\foreach \i in {0,...,2} {
    \foreach \j in {0,...,2} {
        \draw[carnelian!60, fill=carnelian!15] (\j*0.5+6.5,\i*0.5+0.5) rectangle (\j*0.5+6.95,\i*0.5+0.95);
    }
}
\draw[carnelian, fill=carnelian!40, thick] (7.0,1.0) rectangle (7.45,1.45);
\node[font=\small] at (7.25,0.2) {output};
\draw[->, thick, thwsgrey] (3.0,1.25) -- (3.4,1.25);
\draw[->, thick, thwsgrey] (5.5,1.25) -- (6.4,1.25);
\end{tikzpicture}

\end{column}
\begin{column}{0.5\textwidth}


\begin{itemize}
    \item $(i,j)$ --- output position
    \item $(p,q)$ --- index within filter
    \item Two sums: over rows and columns of filter
\end{itemize}
\end{column}
\end{columns}

\vskip 1em

\centering
\structure{Output size:} $(H - k + 1) \times (W - k + 1)$


\end{frame}

\begin{frame}{Convolution as Matrix Multiplication}

\vskip 1em

\structure{Flatten input and output} --- convolution becomes matrix multiplication

\begin{columns}[c]
\begin{column}{0.46\textwidth}
\tikzsetnextfilename{conv2d}
\vspace{-3cm}

% conv2d.tex — placeholder
\begin{tikzpicture}[scale=0.7]
% input image
\foreach \i in {0,...,4} {
    \foreach \j in {0,...,4} {
        \draw[thwspetrol!40, fill=thwspetrol!10] (\j*0.5,\i*0.5) rectangle (\j*0.5+0.45,\i*0.5+0.45);
    }
}
% highlighted patch
\foreach \i in {1,...,3} {
    \foreach \j in {1,...,3} {
        \draw[thwsorange, fill=thwsorange!30, thick] (\j*0.5,\i*0.5) rectangle (\j*0.5+0.45,\i*0.5+0.45);
    }
}
\node[font=\small] at (1.1,-0.3) {input};
% filter
\foreach \i in {0,...,2} {
    \foreach \j in {0,...,2} {
        \draw[thwsorange, fill=thwsorange!50, thick] (\j*0.5+3.5,\i*0.5+0.5) rectangle (\j*0.5+3.95,\i*0.5+0.95);
    }
}
\node[font=\small] at (4.25,0.2) {filter};
% output
\foreach \i in {0,...,2} {
    \foreach \j in {0,...,2} {
        \draw[carnelian!60, fill=carnelian!15] (\j*0.5+6.5,\i*0.5+0.5) rectangle (\j*0.5+6.95,\i*0.5+0.95);
    }
}
\draw[carnelian, fill=carnelian!40, thick] (7.0,1.0) rectangle (7.45,1.45);
\node[font=\small] at (7.25,0.2) {output};
\draw[->, thick, thwsgrey] (3.0,1.25) -- (3.4,1.25);
\draw[->, thick, thwsgrey] (5.5,1.25) -- (6.4,1.25);
\end{tikzpicture}

\end{column}
\begin{column}{0.46\textwidth}
\tikzsetnextfilename{conv2d_matrix}
% conv2d_matrix.tex — full 9x25 weight matrix with 25x1 input and 9x1 output
% kernel values (row-major, bottom-up as in conv2d.tex):
% row0(bot): 1  0 -2
% row1(mid): 3  3 -1
% row2(top): 0  2  0
% kernel flattened row-major (top-to-bottom): 0,2,0,3,3,-1,1,0,-2
% for output position (i,j), active input cols in 25-element flat vector:
% flat index = row*5 + col
% output row k (k=0..8): i=k/3, j=k%3
% active flat indices: (i+0)*5+(j+0), (i+0)*5+(j+1), (i+0)*5+(j+2),
%                      (i+1)*5+(j+0), (i+1)*5+(j+1), (i+1)*5+(j+2),
%                      (i+2)*5+(j+0), (i+2)*5+(j+1), (i+2)*5+(j+2)

\begin{tikzpicture}[scale=0.75]

\def\cw{0.30}  % cell width
\def\ch{0.30}  % cell height
\def\gap{0.15} % gap between blocks

% kernel values in flat order (top-left to bottom-right of filter)
% kv[0..8]: positions in filter scanned row by row
% filter from conv2d.tex (top row first):
% top:    0  2  0
% mid:    3  3 -1
% bot:    1  0 -2
\def\kva{0}
\def\kvb{2}
\def\kvc{0}
\def\kvd{3}
\def\kve{3}
\def\kvf{-1}
\def\kvg{1}
\def\kvh{0}
\def\kvi{-2}

% active flat input indices for each of 9 output positions
% output k: i=floor(k/3), j=mod(k,3)
% active: (i)*5+j, i*5+j+1, i*5+j+2, (i+1)*5+j, ..., (i+2)*5+j+2

% ─────────────────────────────────────────────────────
% WEIGHT MATRIX  9 rows x 25 cols
% ─────────────────────────────────────────────────────
\foreach \outrow in {0,...,8} {
    % compute patch start: i=floor(outrow/3), j=mod(outrow,3)
    \pgfmathtruncatemacro\pi{int(\outrow/3)}
    \pgfmathtruncatemacro\pj{int(mod(\outrow,3))}
    % 9 active flat indices
    \pgfmathtruncatemacro\aa{int(\pi*5+\pj)}
    \pgfmathtruncatemacro\ab{int(\pi*5+\pj+1)}
    \pgfmathtruncatemacro\ac{int(\pi*5+\pj+2)}
    \pgfmathtruncatemacro\ba{int((\pi+1)*5+\pj)}
    \pgfmathtruncatemacro\bb{int((\pi+1)*5+\pj+1)}
    \pgfmathtruncatemacro\bc{int((\pi+1)*5+\pj+2)}
    \pgfmathtruncatemacro\ca{int((\pi+2)*5+\pj)}
    \pgfmathtruncatemacro\cb{int((\pi+2)*5+\pj+1)}
    \pgfmathtruncatemacro\cc{int((\pi+2)*5+\pj+2)}
    \foreach \incol in {0,...,24} {
        % check if active
        \pgfmathtruncatemacro\act{ifthenelse(
            \incol==\aa || \incol==\ab || \incol==\ac ||
            \incol==\ba || \incol==\bb || \incol==\bc ||
            \incol==\ca || \incol==\cb || \incol==\cc, 1, 0)}
        \ifnum\act=1
            \draw[thwsorange, fill=thwsorange!35, thick]
                (\incol*\cw, -\outrow*\ch)
                rectangle (\incol*\cw+\cw-0.02, -\outrow*\ch+\ch-0.02);
            % kernel value: which position in 3x3 kernel?
            \pgfmathtruncatemacro\kidx{ifthenelse(\incol==\aa,0,
                ifthenelse(\incol==\ab,1,
                ifthenelse(\incol==\ac,2,
                ifthenelse(\incol==\ba,3,
                ifthenelse(\incol==\bb,4,
                ifthenelse(\incol==\bc,5,
                ifthenelse(\incol==\ca,6,
                ifthenelse(\incol==\cb,7,8))))))))}
            \ifnum\kidx=0 \def\kval{\kva} \fi
            \ifnum\kidx=1 \def\kval{\kvb} \fi
            \ifnum\kidx=2 \def\kval{\kvc} \fi
            \ifnum\kidx=3 \def\kval{\kvd} \fi
            \ifnum\kidx=4 \def\kval{\kve} \fi
            \ifnum\kidx=5 \def\kval{\kvf} \fi
            \ifnum\kidx=6 \def\kval{\kvg} \fi
            \ifnum\kidx=7 \def\kval{\kvh} \fi
            \ifnum\kidx=8 \def\kval{\kvi} \fi
            \node[font=\tiny] at
                (\incol*\cw+\cw/2-0.01, -\outrow*\ch+\ch/2-0.01) {\kval};
        \else
            \draw[thwsgrey!25, fill=thwsgrey!8]
                (\incol*\cw, -\outrow*\ch)
                rectangle (\incol*\cw+\cw-0.02, -\outrow*\ch+\ch-0.02);
            \node[font=\tiny, thwsgrey!60] at
                (\incol*\cw+\cw/2-0.01, -\outrow*\ch+\ch/2-0.01) {0};
        \fi
    }
}
\node[font=\scriptsize, thwsorange] at (25*\cw/2, 0.45) {$\mW$ \quad $9\times25$};

% ─────────────────────────────────────────────────────
% INPUT VECTOR  25x1  (to the right of matrix)
% ─────────────────────────────────────────────────────
\pgfmathsetmacro\xinput{25*\cw + \gap}
% patch for output 0: active flat indices 0,1,2,5,6,7,10,11,12
\foreach \i in {0,...,24} {
    \pgfmathtruncatemacro\act{ifthenelse(
        \i==0||\i==1||\i==2||
        \i==5||\i==6||\i==7||
        \i==10||\i==11||\i==12, 1, 0)}
    \ifnum\act=1
        \draw[thwspetrol, fill=thwspetrol!35, thick]
            (\xinput, -\i*\cw) rectangle (\xinput+\cw-0.02, -\i*\cw+\cw-0.02);
    \else
        \draw[thwspetrol!25, fill=thwspetrol!8]
            (\xinput, -\i*\cw) rectangle (\xinput+\cw-0.02, -\i*\cw+\cw-0.02);
    \fi
}
\node[font=\scriptsize, thwspetrol] at (\xinput+\cw/2, 0.45) {$\vx$};
\node[font=\scriptsize, thwspetrol] at (\xinput+\cw/2, -25*\cw-0.2) {$25\times1$};

% multiplication sign
\node[font=\normalsize] at (\xinput-\gap/2, -4*\cw) {$\times$};

% ─────────────────────────────────────────────────────
% OUTPUT VECTOR  9x1  (to the left of matrix)
% ─────────────────────────────────────────────────────
\pgfmathsetmacro\xoutput{-\cw - \gap}
\foreach \i in {0,...,8} {
    \ifnum\i=0
        \draw[carnelian, fill=carnelian!40, thick]
            (\xoutput, -\i*\ch) rectangle (\xoutput+\cw-0.02, -\i*\ch+\ch-0.02);
    \else
        \draw[carnelian!35, fill=carnelian!10]
            (\xoutput, -\i*\ch) rectangle (\xoutput+\cw-0.02, -\i*\ch+\ch-0.02);
    \fi
}
\node[font=\scriptsize, carnelian] at (\xoutput+\cw/2, 0.45) {$\vh$};
\node[font=\scriptsize, carnelian] at (\xoutput+\cw/2, -9*\ch-0.2) {$9\times1$};

% equals sign
\node[font=\normalsize] at (\xoutput+\cw+\gap/2, -4*\cw) {$=$};

% ─────────────────────────────────────────────────────
% ANNOTATIONS
% ─────────────────────────────────────────────────────
\draw[decorate, decoration={brace,amplitude=4pt}, thwsorange]
    (0, 0.35) -- (25*\cw, 0.35);
\node[font=\scriptsize, thwsorange, align=center] at (25*\cw/2, 0.75)
    {same kernel values, shifted --- \textbf{tied weights}};

\draw[decorate, decoration={brace,amplitude=4pt,mirror}, thwsgrey]
    (0, -9*\ch-0.05) -- (25*\cw, -9*\ch-0.05);
\node[font=\scriptsize, thwsgrey] at (25*\cw/2, -9*\ch-0.45)
    {most entries zero --- \textbf{sparse}};

\end{tikzpicture}

\end{column}
\end{columns}

\vspace{-3cm}%
\begin{minipage}{1\textwidth}
\centering
Convolution $=$ FC layer with \alert{sparse} and \alert{tied} weights\\[0.5em]
Convolution - \alert{linear transformation!}
\end{minipage}

\end{frame}

\begin{frame}{Padding}

\vskip 1em

\structure{Problem:} convolution shrinks spatial dimensions --- each layer loses $k-1$ pixels

\vskip 1em
\centering
\resizebox{0.6\textwidth}{!}{%
    \tikzsetnextfilename{conv2d_steps}
    % conv2d_steps.tex — three filter positions on same input
\begin{tikzpicture}[scale=0.72]

% --- position 1: top-left (patch at col 0-2, row 2-4 in bottom-up coords) ---
\begin{scope}[xshift=0cm]
    % full grid
    \foreach \col in {0,...,4} {
        \foreach \row in {0,...,4} {
            \draw[thwspetrol!25, fill=thwspetrol!7]
                (\col*0.52, \row*0.52) rectangle (\col*0.52+0.48, \row*0.52+0.48);
        }
    }
    % patch cols 0-2, rows 2-4 (bottom-up: rows 0-2)
    \foreach \col in {0,1,2} {
        \foreach \row in {2,3,4} {
            \draw[thwsorange, fill=thwsorange!30, thick]
                (\col*0.52, \row*0.52) rectangle (\col*0.52+0.48, \row*0.52+0.48);
        }
    }
    % output grid 3x3
    \foreach \col in {0,1,2} {
        \foreach \row in {0,1,2} {
            \draw[carnelian!30, fill=carnelian!8]
                (\col*0.52+2.9, \row*0.52+0.52) rectangle
                (\col*0.52+3.38, \row*0.52+1.0);
        }
    }
    % highlight top-left output
    \draw[carnelian, fill=carnelian!40, thick]
        (2.9, 1.54) rectangle (3.38, 2.02);
    \node[font=\scriptsize, thwsorange] at (2.7,-0.3) {step 1};

\draw[->, thick, thwsgrey] (4.7,1.3) -- (5.5,1.3);

\end{scope}

% --- position 2: shifted right (patch at col 1-3, row 2-4) ---
\begin{scope}[xshift=6cm]
    \foreach \col in {0,...,4} {
        \foreach \row in {0,...,4} {
            \draw[thwspetrol!25, fill=thwspetrol!7]
                (\col*0.52, \row*0.52) rectangle (\col*0.52+0.48, \row*0.52+0.48);
        }
    }
    \foreach \col in {1,2,3} {
        \foreach \row in {2,3,4} {
            \draw[thwsorange, fill=thwsorange!30, thick]
                (\col*0.52, \row*0.52) rectangle (\col*0.52+0.48, \row*0.52+0.48);
        }
    }
    \foreach \col in {0,1,2} {
        \foreach \row in {0,1,2} {
            \draw[carnelian!30, fill=carnelian!8]
                (\col*0.52+2.9, \row*0.52+0.52) rectangle
                (\col*0.52+3.38, \row*0.52+1.0);
        }
    }
    % highlight second output in top row
    \draw[carnelian, fill=carnelian!40, thick]
        (3.42, 1.54) rectangle (3.9, 2.02);
    \node[font=\scriptsize, thwsorange] at (2.7,-0.3) {step 2};

\draw[->, thick, thwsgrey] (4.7,1.3) -- (5.5,1.3);
\end{scope}

% --- position 3: shifted down (patch at col 0-2, row 1-3) ---
\begin{scope}[xshift=12cm]
    \foreach \col in {0,...,4} {
        \foreach \row in {0,...,4} {
            \draw[thwspetrol!25, fill=thwspetrol!7]
                (\col*0.52, \row*0.52) rectangle (\col*0.52+0.48, \row*0.52+0.48);
        }
    }
    \foreach \col in {0,1,2} {
        \foreach \row in {1,2,3} {
            \draw[thwsorange, fill=thwsorange!30, thick]
                (\col*0.52, \row*0.52) rectangle (\col*0.52+0.48, \row*0.52+0.48);
        }
    }
    \foreach \col in {0,1,2} {
        \foreach \row in {0,1,2} {
            \draw[carnelian!30, fill=carnelian!8]
                (\col*0.52+2.9, \row*0.52+0.52) rectangle
                (\col*0.52+3.38, \row*0.52+1.0);
        }
    }
    % highlight second output in first column
    \draw[carnelian, fill=carnelian!40, thick]
        (2.9, 1.04) rectangle (3.38, 1.52);
    \node[font=\scriptsize, thwsorange] at (2.7,-0.3) {step 3};
\end{scope}

% \node[font=\scriptsize, thwspetrol, align=center] at (8.5,-1)
%     {same filter slides across input --- each position produces one output value};

\end{tikzpicture}
%
}

\vskip 1em


\begin{columns}[t]
\begin{column}{0.5\textwidth}
\tikzsetnextfilename{conv2d_padding}
% conv2d_padding.tex
% Left: no padding, 3x3 input, 3x3 filter -> 1x1 output
% Right: same padding (1), 3x3 input + zeros -> 5x5, 3x3 filter -> 3x3 output
\begin{tikzpicture}[scale=0.72]

% ── LEFT: no padding ──────────────────────────────────────────────
\begin{scope}[xshift=0cm]
    % 3x3 input
    \foreach \col in {0,1,2} {
        \foreach \row in {0,1,2} {
            \draw[thwspetrol, fill=thwspetrol!35, thick]
                (\col*0.52, \row*0.52) rectangle (\col*0.52+0.48, \row*0.52+0.48);
        }
    }
    % filter covers entire input
    \foreach \col in {0,1,2} {
        \foreach \row in {0,1,2} {
            \draw[thwsorange, fill=thwsorange!30, thick]
                (\col*0.52+0.06, \row*0.52+0.06) rectangle
                (\col*0.52+0.44, \row*0.52+0.44);
        }
    }
    % 1x1 output
    \draw[carnelian, fill=carnelian!40, thick]
        (2.0, 0.76) rectangle (2.48, 1.24);
    % labels
    \node[font=\scriptsize, thwspetrol] at (0.76, -0.3) {input $3\times3$};
    \node[font=\scriptsize, carnelian]  at (2.24, -0.3) {out $1\times1$};
    \node[font=\small, thwsgrey] at (0.76, 2.1) {no padding};
    % output size annotation
    \draw[->, thwsgrey, thin] (1.72, 1.0) -- (1.95, 1.0);
\end{scope}

% ── RIGHT: same padding (p=1) ────────────────────────────────────
\begin{scope}[xshift=4.5cm]
    % zero padding ring (5x5)
    \foreach \col in {0,...,4} {
        \foreach \row in {0,...,4} {
            \draw[thwsgrey!30, fill=thwsgrey!10]
                (\col*0.52, \row*0.52) rectangle (\col*0.52+0.48, \row*0.52+0.48);
            \node[font=\tiny, thwsgrey!60] at (\col*0.52+0.24, \row*0.52+0.24) {0};
        }
    }
    % original 3x3 input (center)
    \foreach \col in {1,2,3} {
        \foreach \row in {1,2,3} {
            \draw[thwspetrol, fill=thwspetrol!35, thick]
                (\col*0.52, \row*0.52) rectangle (\col*0.52+0.48, \row*0.52+0.48);
        }
    }
    % filter at top-left position (cols 0-2, rows 2-4)
    \foreach \col in {0,1,2} {
        \foreach \row in {2,3,4} {
            \draw[thwsorange, fill=thwsorange!30, thick]
                (\col*0.52+0.04, \row*0.52+0.04) rectangle
                (\col*0.52+0.46, \row*0.52+0.46);
        }
    }
    % 3x3 output
    \foreach \col in {0,1,2} {
        \foreach \row in {0,1,2} {
            \draw[carnelian!30, fill=carnelian!8]
                (\col*0.52+3.2, \row*0.52+0.52) rectangle
                (\col*0.52+3.68, \row*0.52+1.0);
        }
    }
    \draw[carnelian, fill=carnelian!40, thick]
        (3.2, 2.08) rectangle (3.68, 2.56);
    \draw[->, thwsgrey, thin] (2.7, 1.8) -- (3.15, 1.8);
    % labels
    \node[font=\scriptsize, thwsgrey]   at (1.2, -0.3) {padded $5\times5$};
    \node[font=\scriptsize, carnelian]  at (3.7, -0.3) {out $3\times3$};
    \node[font=\small, thwsgrey] at (1.2, 2.8) {same padding ($p=1$)};
\end{scope}

\end{tikzpicture}

\end{column}
\begin{column}{0.5\textwidth}
\structure{Valid padding} ($p=0$)
\begin{itemize}
    \item No zeros added
    \item Output shrinks: $(H-k+1) \times (W-k+1)$
    \item Spatial information lost at borders
\end{itemize}
\structure{Same padding} ($p=(k-1)/2$)
\begin{itemize}
    \item Zeros added around border
    \item Output same size as input
    \item Standard choice in most architectures
\end{itemize}
\end{column}
\end{columns}

\end{frame}

\begin{frame}{Stride}

\vskip 1em

\structure{Stride $s$:} filter $\kernel$ moves $s$ pixels at each step instead of 1

\vskip 1em

\tikzsetnextfilename{conv2d_stride}
% conv2d_stride.tex
% Top row: stride=1, three consecutive positions
% Bottom row: stride=2, two positions (jumps over one)
\begin{tikzpicture}[scale=0.72]

% ── TOP: stride=1 ────────────────────────────────────────────────
\node[font=\small, thwspetrol] at (-0.4, 2.8) {$s{=}1$};

\foreach \step/\xsh/\pcol in {0/0/0, 1/5/1, 2/10/2} {
    \begin{scope}[xshift=\xsh cm]
        \foreach \col in {0,...,4} {
            \foreach \row in {0,...,4} {
                \draw[thwspetrol!25, fill=thwspetrol!7]
                    (\col*0.52, \row*0.52) rectangle (\col*0.52+0.48, \row*0.52+0.48);
            }
        }
        % patch starts at col \pcol
        \foreach \col in {0,1,2} {
            \foreach \row in {2,3,4} {
                \pgfmathtruncatemacro\dc{\col+\pcol}
                \draw[thwsorange, fill=thwsorange!30, thick]
                    (\dc*0.52, \row*0.52) rectangle (\dc*0.52+0.48, \row*0.52+0.48);
            }
        }
        % output grid
        \foreach \col in {0,1,2} {
            \foreach \row in {0,1,2} {
                \draw[carnelian!30, fill=carnelian!8]
                    (\col*0.52+3.0, \row*0.52+0.52) rectangle
                    (\col*0.52+3.48, \row*0.52+1.0);
            }
        }
        % highlight active output cell
        \pgfmathsetmacro\ox{3.0+\pcol*0.52}
        \draw[carnelian, fill=carnelian!40, thick]
            (\ox, 1.56) rectangle (\ox+0.48, 2.04);
        \draw[->, thwsgrey, thin] (2.7,1.3) -- (2.95,1.3);
        % step label
        \node[font=\scriptsize, thwsorange] at (1.2,-0.3) {step \pgfmathparse{int(\step+1)}\pgfmathresult};
    \end{scope}
    % arrow between steps
    \ifnum\step<2
        \pgfmathsetmacro\ax{\xsh+4.5}
        \draw[->, thick, thwsgrey] (\ax+0.1,1.3) -- (\ax+0.4,1.3);
    \fi
}

% ── BOTTOM: stride=2 ─────────────────────────────────────────────
\node[font=\small, thwspetrol] at (-0.4, -1.0) {$s{=}2$};

\foreach \step/\xsh/\pcol in {0/0/0, 1/5/2} {
    \begin{scope}[xshift=\xsh cm, yshift=-4.0cm]
        \foreach \col in {0,...,4} {
            \foreach \row in {0,...,4} {
                \draw[thwspetrol!25, fill=thwspetrol!7]
                    (\col*0.52, \row*0.52) rectangle (\col*0.52+0.48, \row*0.52+0.48);
            }
        }
        \foreach \col in {0,1,2} {
            \foreach \row in {2,3,4} {
                \pgfmathtruncatemacro\dc{\col+\pcol}
                \draw[thwsorange, fill=thwsorange!30, thick]
                    (\dc*0.52, \row*0.52) rectangle (\dc*0.52+0.48, \row*0.52+0.48);
            }
        }
        % output grid (stride=2: 2x2 output for 5x5 input, k=3)
        \foreach \col in {0,1} {
            \foreach \row in {0,1} {
                \draw[carnelian!30, fill=carnelian!8]
                    (\col*0.52+3.0, \row*0.52+0.78) rectangle
                    (\col*0.52+3.48, \row*0.52+1.26);
            }
        }
        % highlight active output
        \pgfmathtruncatemacro\oc{int(\pcol/2)}
        \pgfmathsetmacro\ox{3.0+\oc*0.52}
        \draw[carnelian, fill=carnelian!40, thick]
            (\ox, 1.3) rectangle (\ox+0.48, 1.78);
        \draw[->, thwsgrey, thin] (2.7,1.3) -- (2.95,1.3);
        \node[font=\scriptsize, thwsorange] at (1.2,-0.3) {step \pgfmathparse{int(\step+1)}\pgfmathresult};
    \end{scope}
    \ifnum\step<1
        \pgfmathsetmacro\ax{\xsh+3.6}
        \draw[->, thick, thwsgrey] (\ax+1.0,-2.7) -- (\ax+1.3,-2.7);
        \node[font=\scriptsize, thwsorange] at (\ax+2.3,-1.2) {skip 1 pixel};
    \fi
}

% output size annotations
\node[font=\scriptsize, carnelian] at (13.9, 2.3) {output $3\times3$};
\node[font=\scriptsize, carnelian] at (8.7, -1.9) {output $2\times2$};

\end{tikzpicture}


\vspace{-2.5cm}\hspace{8cm}%
\begin{minipage}{0.4\textwidth}
\structure{Stride 1:} $\kernel$ moves one step a time\\[0.3em]
\structure{Stride 2:} $\kernel$ skips one - output halved\\[0.5em]
\end{minipage}

\vspace{0.5cm}%
\[
   \hspace{6cm} \text{\alert{\textbf{General output size:}}} \quad \left\lfloor \frac{H + 2p - k}{s} \right\rfloor + 1
\]

\end{frame}

\begin{frame}{Multiple Input Channels}

\vskip 1em

\structure{RGB input:} $C_{in}$ channels --- filter must have same depth $C_{in}$

\vskip 1em

\structure{Operation:} filter has $C_{in}$ slices --- each slice convolves with one input channel

\[
    (\mX * \kernel)[i,j] = \sum_{c=0}^{C_{in}-1}\underbrace{\sum_{p=0}^{k-1}\sum_{q=0}^{k-1}
    \mX[c,\, i+p,\, j+q]\; \kernel[c,\, p,\, q]}_{\text{conv with channel } c}
\]

\begin{itemize}
    \item Inner double sum: convolution of filter slice $c$ with input channel $c$
    \item Outer sum over $c$: results summed across all channels
    \item One filter $\to$ one output feature map
\end{itemize}

\begin{center}
\tikzsetnextfilename{conv2d_channels}
% conv2d_channels.tex
% Shows single filter with C_in=3 channels -> one output feature map
% Input: 3 stacked 5x5 grids (petrol shades)
% Filter: 3 stacked 3x3 grids (orange)
% Output: one 3x3 feature map (carnelian)
\begin{tikzpicture}[scale=0.72]

\def\cs{0.42}  % cell size
\def\os{0.12}  % channel offset (3D effect)

% ── INPUT VOLUME: 3 channels ──────────────────────────────────────
% draw back to front so front channel appears on top
\foreach \ch in {2,1,0} {
    \pgfmathsetmacro\xoff{\ch*\os}
    \pgfmathsetmacro\yoff{\ch*\os}
    \pgfmathtruncatemacro\shade{15+\ch*10}
    \foreach \col in {0,...,4} {
        \foreach \row in {0,...,4} {
            \draw[thwspetrol!\shade!black, fill=thwspetrol!\shade]
                (\col*\cs+\xoff, \row*\cs+\yoff)
                rectangle (\col*\cs+\cs-0.02+\xoff, \row*\cs+\cs-0.02+\yoff);
        }
    }
}
\node[font=\scriptsize, thwspetrol] at (1.1+2*\os, -0.35) {input $C_{in}{\times}H{\times}W$};

% ── ARROW ────────────────────────────────────────────────────────
\draw[->, thick, thwsgrey] (2.6, 1.2) -- (3.1, 1.2);

% ── FILTER: 3 layers (one per input channel) ─────────────────────
\foreach \ch in {2,1,0} {
    \pgfmathsetmacro\xoff{3.3+\ch*\os}
    \pgfmathsetmacro\yoff{0.5+\ch*\os}
    \foreach \col in {0,1,2} {
        \foreach \row in {0,1,2} {
            \draw[thwsorange, fill=thwsorange!30, thick]
                (\col*\cs+\xoff, \row*\cs+\yoff)
                rectangle (\col*\cs+\cs-0.02+\xoff, \row*\cs+\cs-0.02+\yoff);
        }
    }
}
\node[font=\scriptsize, thwsorange] at (4.25, -0.35) {filter $C_{in}{\times}k{\times}k$};

% ── ARROW ────────────────────────────────────────────────────────
\draw[->, thick, thwsgrey] (5.7, 1.2) -- (6.2, 1.2);

% ── OUTPUT: one feature map ───────────────────────────────────────
\foreach \col in {0,1,2} {
    \foreach \row in {0,1,2} {
        \draw[carnelian, fill=carnelian!25]
            (\col*\cs+6.4, \row*\cs+0.5)
            rectangle (\col*\cs+\cs-0.02+6.4, \row*\cs+\cs-0.02+0.5);
    }
}
\node[font=\scriptsize, carnelian] at (7.1, -0.35) {feature map $H'{\times}W'$};


\end{tikzpicture}

\end{center}

\end{frame}

\begin{frame}{Multiple Output Channels}

\vskip 1em

\structure{$C_{out}$ filters:} each filter independently applied to the full input volume

\vskip 0.5em

\[
    \vh_f[i,j] = \sum_{c=0}^{C_{in}-1}\sum_{p=0}^{k-1}\sum_{q=0}^{k-1}
    \mX[c,\, i+p,\, j+q]\; \kernel_f[c,\, p,\, q]
    \quad f = 0,\ldots,C_{out}-1
\]

\begin{center}
\tikzsetnextfilename{conv2d_filters}
% conv2d_filters.tex
% Shows C_out filters -> C_out feature maps
% Input: 3 stacked grids (petrol)
% Filters: C_out=3 filter volumes stacked with offset (orange)
% Output: C_out=3 feature maps stacked (carnelian)
\begin{tikzpicture}[scale=0.72]

\def\cs{0.42}   % cell size
\def\os{0.12}   % channel offset
\def\fs{0.10}   % filter stack offset

% ── INPUT VOLUME ──────────────────────────────────────────────────
\foreach \ch in {2,1,0} {
    \pgfmathsetmacro\xoff{\ch*\os}
    \pgfmathsetmacro\yoff{\ch*\os}
    \pgfmathtruncatemacro\shade{15+\ch*10}
    \foreach \col in {0,...,4} {
        \foreach \row in {0,...,4} {
            \draw[thwspetrol!\shade!black, fill=thwspetrol!\shade]
                (\col*\cs+\xoff, \row*\cs+\yoff)
                rectangle (\col*\cs+\cs-0.02+\xoff, \row*\cs+\cs-0.02+\yoff);
        }
    }
}
\node[font=\scriptsize, thwspetrol] at (1.1+2*\os, -0.35) {input $C_{in}{\times}H{\times}W$};

% ── ARROW ────────────────────────────────────────────────────────
\draw[->, thick, thwsgrey] (2.6, 1.4) -- (3.1, 1.4);

% ── C_out FILTERS stacked ─────────────────────────────────────────
% 3 filter volumes, each offset slightly
\foreach \filt in {2,1,0} {
    \pgfmathsetmacro\xfoff{3.3+\filt*\fs*3}
    \pgfmathsetmacro\yfoff{0.2+\filt*0.45}
    % each filter volume has C_in=3 layers
    \foreach \ch in {2,1,0} {
        \pgfmathsetmacro\xoff{\xfoff+\ch*\os}
        \pgfmathsetmacro\yoff{\yfoff+\ch*\os}
        \foreach \col in {0,1,2} {
            \foreach \row in {0,1,2} {
                \draw[thwsorange, fill=thwsorange!25, thick]
                    (\col*\cs+\xoff, \row*\cs+\yoff)
                    rectangle (\col*\cs+\cs-0.02+\xoff, \row*\cs+\cs-0.02+\yoff);
            }
        }
    }
}
\node[font=\scriptsize, thwsorange] at (4.5, -0.35) {$C_{out}$ filters};

% ── ARROWS ───────────────────────────────────────────────────────
\draw[->, thick, thwsgrey] (5.8, 1.4) -- (6.3, 1.4);

% ── C_out FEATURE MAPS stacked ────────────────────────────────────
\foreach \filt in {2,1,0} {
    \pgfmathsetmacro\xoff{6.5+\filt*\os}
    \pgfmathsetmacro\yoff{0.3+\filt*0.45}
    \foreach \col in {0,1,2} {
        \foreach \row in {0,1,2} {
            \draw[carnelian, fill=carnelian!25]
                (\col*\cs+\xoff, \row*\cs+\yoff)
                rectangle (\col*\cs+\cs-0.02+\xoff, \row*\cs+\cs-0.02+\yoff);
        }
    }
}
\node[font=\scriptsize, carnelian] at (7.4, -0.35) {$C_{out}$ feature maps};

% ── LABEL ────────────────────────────────────────────────────────
\node[font=\scriptsize, thwsgrey, align=center] at (4.2, -0.9)
    {each filter produces one feature map $\to$ output $C_{out}{\times}H'{\times}W'$};

\end{tikzpicture}

\end{center}

\begin{itemize}
    \item Each filter $\kernel_f$ produces one feature map $\vh_f$ independently
    \item Output tensor: $C_{out} \times H' \times W'$
    \item Parameters: $C_{out} \times C_{in} \times k \times k + C_{out}$ (with bias)
\end{itemize}

\end{frame}

\begin{frame}{Receptive Field}

\vskip 2em

\structure{Receptive field:} region of the input that influences a given neuron's output

\begin{center}
\tikzsetnextfilename{receptive_field}
% receptive_field.tex
% Three 7x7 grids showing growing receptive field
% Left: input, 5x5 center highlighted, top-left 3x3 stronger
% Middle: layer1, 3x3 center highlighted, top-left cell stronger
% Right: layer2, center pixel highlighted
% Dashed lines connect highlighted regions between grids

\begin{tikzpicture}[scale=0.72]

\def\cs{0.40}   % cell size
\def\gap{1.2}   % gap between grids

% ── LEFT GRID: input ─────────────────────────────────────────────
\def\xl{0}
\foreach \col in {0,...,6} {
    \foreach \row in {0,...,6} {
        \draw[thwspetrol!20, fill=thwspetrol!6]
            (\col*\cs+\xl, \row*\cs) rectangle
            (\col*\cs+\cs-0.02+\xl, \row*\cs+\cs-0.02);
    }
}
% 5x5 center highlight (cols 1-5, rows 1-5)
\foreach \col in {1,2,3,4,5} {
    \foreach \row in {1,2,3,4,5} {
        \draw[carnelian, fill=carnelian!20, thick]
            (\col*\cs+\xl, \row*\cs) rectangle
            (\col*\cs+\cs-0.02+\xl, \row*\cs+\cs-0.02);
    }
}
% top-left 3x3 stronger (cols 1-3, rows 3-5)
\foreach \col in {1,2,3} {
    \foreach \row in {3,4,5} {
        \draw[thwspetrol, fill=thwspetrol!25, thick]
            (\col*\cs+\xl, \row*\cs) rectangle
            (\col*\cs+\cs-0.02+\xl, \row*\cs+\cs-0.02);
    }
}
\node[font=\scriptsize, thwspetrol] at (\xl+3*\cs, -0.35) {input};

% ── MIDDLE GRID: layer1 ──────────────────────────────────────────
\pgfmathsetmacro\xm{7*\cs+\gap}
\foreach \col in {0,...,6} {
    \foreach \row in {0,...,6} {
        \draw[thwspetrol!20, fill=thwspetrol!6]
            (\col*\cs+\xm, \row*\cs) rectangle
            (\col*\cs+\cs-0.02+\xm, \row*\cs+\cs-0.02);
    }
}
% 3x3 center highlight (cols 2-4, rows 2-4)
\foreach \col in {2,3,4} {
    \foreach \row in {2,3,4} {
        \draw[carnelian, fill=carnelian!20, thick]
            (\col*\cs+\xm, \row*\cs) rectangle
            (\col*\cs+\cs-0.02+\xm, \row*\cs+\cs-0.02);
    }
}
% top-left cell stronger (col=2, row=4)
\draw[thwspetrol, fill=thwspetrol!70, thick]
    (2*\cs+\xm, 4*\cs) rectangle (2*\cs+\cs-0.02+\xm, 4*\cs+\cs-0.02);
\node[font=\scriptsize, thwspetrol] at (\xm+3*\cs, -0.35) {layer 1};
\node[font=\scriptsize, thwspetrol]   at (\xm+3*\cs, -0.65) {RF: $3{\times}3$};

% ── RIGHT GRID: layer2 ───────────────────────────────────────────
\pgfmathsetmacro\xr{\xm+7*\cs+\gap}
\foreach \col in {0,...,6} {
    \foreach \row in {0,...,6} {
        \draw[thwspetrol!20, fill=thwspetrol!6]
            (\col*\cs+\xr, \row*\cs) rectangle
            (\col*\cs+\cs-0.02+\xr, \row*\cs+\cs-0.02);
    }
}
% center pixel
\draw[carnelian, fill=carnelian!60, thick]
    (3*\cs+\xr, 3*\cs) rectangle (3*\cs+\cs-0.02+\xr, 3*\cs+\cs-0.02);
\node[font=\scriptsize, carnelian] at (\xr+3*\cs, -0.35) {layer 2};
\node[font=\scriptsize, carnelian]  at (\xr+3*\cs, -0.65) {RF: $7{\times}7$};

% ── ARROWS between grids ─────────────────────────────────────────
\draw[->, thick, thwsgrey] (7*\cs+0.1, 3*\cs) -- (7*\cs+\gap-0.1, 3*\cs);
\draw[->, thick, thwsgrey] (\xm+7*\cs+0.1, 3*\cs) -- (\xm+7*\cs+\gap-0.1, 3*\cs);

% ── DASHED LINES: right grid center pixel -> middle grid 3x3 ─────
% top-left corner of center pixel in right grid -> top-left of 3x3 in middle
\draw[carnelian, dashed, thin]
    (3*\cs+\xr, 3*\cs+\cs-0.02)    % top-left of center pixel
    -- (2*\cs+\xm+\cs-0.02+7*\cs+\gap-7*\cs-\gap, 4*\cs+\cs-0.02);  % top-left of 3x3
% simpler: compute directly
\draw[carnelian, dashed, thin]
    (3*\cs+\xr, 4*\cs-0.02)         % top-left corner of center pixel
    -- (2*\cs+\xm, 5*\cs-0.02);     % top-left corner of 3x3 highlight
\draw[carnelian, dashed, thin]
    (4*\cs-0.02+\xr, 3*\cs)         % bottom-right corner of center pixel
    -- (5*\cs-0.02+\xm, 2*\cs);     % bottom-right corner of 3x3 highlight

% ── DASHED LINES: middle grid 3x3 -> left grid 5x5 ───────────────
\draw[thwspetrol, dashed, thin]
    (2*\cs+\xm, 5*\cs-0.02)         % top-left of 3x3 in middle
    -- (1*\cs+\xl, 6*\cs-0.02);     % top-left of 5x5 in left
\draw[thwspetrol, dashed, thin]
    (2*\cs+\cs+\xm, 4*\cs-0.02)         % bottom-right of 1x1 in middle
    -- (3*\cs+\xl, 3*\cs-0.02);     % bottom-right of 3x3 in left
\draw[carnelian, dashed, thin]
    (5*\cs-0.02+\xm, 2*\cs)         % bottom-right of 3x3 in middle
    -- (6*\cs-0.02+\xl, 1*\cs);     % bottom-right of 5x5 in left

\end{tikzpicture}

\end{center}

\vskip 0.5em

\begin{itemize}
    \item Each $3\times3$ conv layer adds 2 pixels to the receptive field in each dimension
    \item $L$ layers of $3\times3$ filters: receptive field $= (2L+1)\times(2L+1)$
    \item Deep networks see large regions despite small local filters
    \item Early layers: local features (edges) --- later layers: global patterns (objects)
\end{itemize}

\end{frame}


% -------------------------------------------------------
\section{Pooling}
% -------------------------------------------------------


\begin{frame}{}
\vskip 8em
\begin{center}
{\Large \structure{Pooling}}
\end{center}
\end{frame}

\begin{frame}{Pooling}

\vskip 2em

\structure{Pooling:} spatial downsampling --- reduces feature map size

\vskip 1em

\begin{columns}[t]
\begin{column}{0.46\textwidth}
\structure{Max pooling:}
\[
    y[i,j] = \max_{p,q \in \mathcal{R}_{ij}} x[p,q]
\]
\begin{itemize}
    \item Takes maximum over local region
    \item Retains strongest activation
    \item Most commonly used
\end{itemize}
\end{column}
\begin{column}{0.46\textwidth}
\structure{Average pooling:}
\[
    y[i,j] = \frac{1}{|\mathcal{R}_{ij}|}\sum_{p,q \in \mathcal{R}_{ij}} x[p,q]
\]
\begin{itemize}
    \item Takes average over local region
    \item Smoother downsampling
    \item Less common in hidden layers
\end{itemize}
\end{column}
\begin{column}{0.04\textwidth}
\end{column}
\end{columns}

\vskip 1.5em

\structure{Benefits:}
\begin{itemize}
    \item Reduces spatial dimensions --- fewer parameters in subsequent layers
    \item Local translation invariance --- small shifts in feature position do not affect output
\end{itemize}

\end{frame}

% % -------------------------------------------------------

% \begin{frame}{Global Average Pooling}

% \vskip 2em

% \structure{Global average pooling (GAP):} averages entire feature map to a single value per channel

% \[
%     y_c = \frac{1}{H \cdot W} \sum_{i=1}^{H}\sum_{j=1}^{W} x_c[i,j]
% \]

% \begin{itemize}
%     \item Output: vector of length $C_{\text{out}}$ --- one value per feature map
%     \item Replaces large FC layers at the end of the network
%     \item Dramatically reduces parameters
%     \item Forces each feature map to represent a single concept
% \end{itemize}

% \vskip 1em

% \structure{Used in:} ResNet, GoogLeNet, MobileNet --- modern standard

% \vskip 0.5em

% \structure{In PyTorch:} \texttt{nn.AdaptiveAvgPool2d((1,1))} followed by \texttt{flatten}

% \end{frame}

% % -------------------------------------------------------
% \section{CNN Architecture}
% % -------------------------------------------------------

% \begin{frame}{CNN Building Block}

% \vskip 2em

% \structure{Standard pattern:} Conv $\to$ BatchNorm $\to$ ReLU $\to$ (Pool)

% \vskip 1em

% \begin{center}
% % cnn_block.tex — placeholder
\begin{tikzpicture}[node distance=1.2cm, every node/.style={font=\small}]
\node[draw, thwspetrol, fill=thwspetrol!15, rounded corners, minimum width=1.5cm, minimum height=0.6cm] (conv) {Conv};
\node[draw, thwsorange, fill=thwsorange!15, rounded corners, minimum width=1.5cm, minimum height=0.6cm, right of=conv] (bn) {BN};
\node[draw, carnelian, fill=carnelian!15, rounded corners, minimum width=1.5cm, minimum height=0.6cm, right of=bn] (relu) {ReLU};
\node[draw, thwsblue, fill=thwsblue!10, rounded corners, minimum width=1.5cm, minimum height=0.6cm, right of=relu] (pool) {Pool};
\draw[->, thick] (conv) -- (bn);
\draw[->, thick] (bn) -- (relu);
\draw[->, thick] (relu) -- (pool);
\draw[->, thick] (-1.5,0) -- (conv);
\draw[->, thick] (pool) -- ++(1.2,0);
\end{tikzpicture}

% \end{center}

% \vskip 1em

% \begin{itemize}
%     \item \structure{Conv:} extract features via learned filters
%     \item \structure{BatchNorm:} stabilise activations --- essential for deep CNNs
%     \item \structure{ReLU:} non-linearity --- enables complex feature combinations
%     \item \structure{Pool:} spatial downsampling --- reduce resolution, increase receptive field
% \end{itemize}

% \end{frame}

% % -------------------------------------------------------

% \begin{frame}{Full CNN Architecture}

% \vskip 2em

% \structure{Two stages:} feature extraction (conv blocks) $\to$ classification (FC head)

% \vskip 1em

% \begin{center}
% % cnn_architecture.tex — placeholder
\begin{tikzpicture}[node distance=1.0cm, every node/.style={font=\scriptsize}]
\node[draw, thwspetrol!60, fill=thwspetrol!10, minimum width=0.8cm, minimum height=1.6cm] (img) {};
\node[below=0.05cm of img] {$3\times32\times32$};
\node[draw, thwspetrol, fill=thwspetrol!20, minimum width=0.6cm, minimum height=1.2cm, right of=img, node distance=1.4cm] (c1) {};
\node[below=0.05cm of c1] {$32\times16\times16$};
\node[draw, thwspetrol, fill=thwspetrol!30, minimum width=0.5cm, minimum height=0.9cm, right of=c1, node distance=1.2cm] (c2) {};
\node[below=0.05cm of c2] {$64\times8\times8$};
\node[draw, thwspetrol, fill=thwspetrol!40, minimum width=0.4cm, minimum height=0.6cm, right of=c2, node distance=1.1cm] (c3) {};
\node[below=0.05cm of c3] {$128\times4\times4$};
\node[draw, carnelian, fill=carnelian!15, minimum width=0.8cm, minimum height=0.4cm, right of=c3, node distance=1.2cm] (fc1) {};
\node[below=0.05cm of fc1] {FC 256};
\node[draw, carnelian, fill=carnelian!25, minimum width=0.6cm, minimum height=0.3cm, right of=fc1, node distance=1.1cm] (fc2) {};
\node[below=0.05cm of fc2] {FC 10};
\draw[->, thick] (img) -- (c1);
\draw[->, thick] (c1) -- (c2);
\draw[->, thick] (c2) -- (c3);
\draw[->, thick] (c3) -- (fc1);
\draw[->, thick] (fc1) -- (fc2);
\node[above=0.3cm of c2, xshift=-0.5cm, font=\scriptsize, thwspetrol] {conv blocks};
\node[above=0.3cm of fc1, font=\scriptsize, carnelian] {FC head};
\end{tikzpicture}

% \end{center}

% \vskip 1em

% \begin{columns}[t]
% \begin{column}{0.46\textwidth}
% \structure{Conv blocks:}
% \begin{itemize}
%     \item Increasing channels: 3 $\to$ 32 $\to$ 64 $\to$ 128
%     \item Decreasing spatial size via pooling
%     \item Learn hierarchical features
% \end{itemize}
% \end{column}
% \begin{column}{0.46\textwidth}
% \structure{FC head:}
% \begin{itemize}
%     \item Flatten spatial features
%     \item One or more FC layers
%     \item Final softmax for classification
% \end{itemize}
% \end{column}
% \begin{column}{0.04\textwidth}
% \end{column}
% \end{columns}

% \end{frame}

% % -------------------------------------------------------
% \section{Classic Architectures}
% % -------------------------------------------------------

% \begin{frame}{LeNet-5 (LeCun et al., 1998)}

% \vskip 2em

% \structure{First successful CNN} --- handwritten digit recognition (MNIST)

% \vskip 1em

% \begin{itemize}
%     \item 2 conv layers + 2 FC layers
%     \item $28\times28$ input $\to$ 10 class output
%     \item $\sim$60k parameters --- tiny by modern standards
%     \item Used sigmoid/tanh activations, average pooling
% \end{itemize}

% \vskip 1em

% \structure{Significance:}
% \begin{itemize}
%     \item Demonstrated that learned features outperform hand-crafted ones
%     \item Introduced the conv $\to$ pool $\to$ conv $\to$ pool $\to$ FC pattern
%     \item Deployed commercially for cheque reading
% \end{itemize}

% \vskip 1em

% \structure{Limitation:} too shallow and small for complex tasks --- needed more compute

% \end{frame}

% % -------------------------------------------------------

% \begin{frame}{AlexNet (Krizhevsky et al., 2012)}

% \vskip 2em

% \structure{ImageNet moment} --- won ILSVRC 2012 with 15.3\% top-5 error vs 26.2\% runner-up

% \vskip 1em

% \begin{itemize}
%     \item 5 conv layers + 3 FC layers, $\sim$60M parameters
%     \item Key innovations:
%     \begin{itemize}
%         \item \alert{ReLU activations} instead of sigmoid/tanh --- faster training
%         \item \alert{Dropout} in FC layers --- regularisation
%         \item \alert{Data augmentation} --- random crops, horizontal flips
%         \item \alert{GPU training} --- two GTX 580 GPUs
%         \item Local response normalisation (since superseded by BatchNorm)
%     \end{itemize}
% \end{itemize}

% \vskip 1em

% \structure{Impact:} launched the modern deep learning era --- proved deep CNNs work at scale

% \end{frame}

% % -------------------------------------------------------

% \begin{frame}{VGG (Simonyan \& Zisserman, 2014)}

% \vskip 2em

% \structure{Simplicity through depth} --- all $3\times3$ convolutions, increasing depth

% \vskip 1em

% \begin{itemize}
%     \item VGG-16: 13 conv layers + 3 FC layers, $\sim$138M parameters
%     \item Only $3\times3$ filters throughout --- simplicity as a design principle
%     \item Depth: 16--19 layers
% \end{itemize}

% \vskip 1em

% \structure{Key insight:} two $3\times3$ conv layers have same receptive field as one $5\times5$ but:
% \begin{itemize}
%     \item Fewer parameters: $2 \times 3^2 = 18$ vs $5^2 = 25$
%     \item More non-linearities --- more expressive
% \end{itemize}

% \vskip 1em

% \structure{Limitation:} 3 large FC layers contain most parameters ($\sim$100M) --- memory intensive

% \vskip 0.5em

% \structure{Legacy:} VGG features widely used for transfer learning; clean design still instructive

% \end{frame}

% % -------------------------------------------------------

% \begin{frame}{ResNet (He et al., 2016)}

% \vskip 2em

% \structure{Problem:} deeper networks were performing \alert{worse} than shallower ones

% \vskip 1em

% \structure{Degradation problem:} adding more layers hurts training accuracy --- not overfitting, but optimisation failure

% \vskip 1em

% \structure{Key idea:} \textbf{residual connections} (skip connections)

% \[
%     \vy = \mathcal{F}(\vx, \{\mW_i\}) + \vx
% \]

% \begin{itemize}
%     \item Instead of learning $\vy = \mathcal{F}(\vx)$, learn the \textbf{residual} $\mathcal{F}(\vx) = \vy - \vx$
%     \item If optimal solution is identity, easier to learn $\mathcal{F}(\vx) = 0$ than $\mathcal{F}(\vx) = \vx$
%     \item Gradient flows directly through skip connection --- no vanishing gradients
% \end{itemize}

% \vskip 0.5em

% \structure{Result:} successfully trained networks with 50, 101, 152 layers

% \end{frame}

% % -------------------------------------------------------

% \begin{frame}{Residual Block}

% \vskip 2em

% \begin{columns}[c]
% \begin{column}{0.55\textwidth}

% \structure{Basic residual block:}
% \[
%     \vy = \text{ReLU}(\mathcal{F}(\vx) + \vx)
% \]

% \vskip 0.5em

% $\mathcal{F}(\vx)$: two Conv-BN-ReLU layers

% \vskip 1em

% \structure{Bottleneck block} (ResNet-50+):
% \begin{itemize}
%     \item $1\times1$ conv: reduce channels
%     \item $3\times3$ conv: spatial processing
%     \item $1\times1$ conv: restore channels
% \end{itemize}
% Fewer parameters than basic block at same depth

% \vskip 1em

% \structure{Projection shortcut:} when dimensions change, use $1\times1$ conv on skip connection to match

% \end{column}
% \begin{column}{0.4\textwidth}
% \centering
% % residual_block.tex — placeholder
\begin{tikzpicture}[node distance=1.0cm, every node/.style={font=\small}]
\node[draw, thwspetrol, fill=thwspetrol!15, rounded corners, minimum width=2cm, minimum height=0.6cm] (conv1) {Conv-BN-ReLU};
\node[draw, thwspetrol, fill=thwspetrol!15, rounded corners, minimum width=2cm, minimum height=0.6cm, below of=conv1, node distance=1.2cm] (conv2) {Conv-BN};
\node[circle, draw, thwsorange, fill=thwsorange!20, minimum size=0.5cm, below of=conv2, node distance=1.0cm] (add) {$+$};
\node[draw, carnelian, fill=carnelian!15, rounded corners, minimum width=2cm, minimum height=0.6cm, below of=add, node distance=0.8cm] (relu) {ReLU};

\draw[->, thick] (conv1) -- (conv2);
\draw[->, thick] (conv2) -- (add);
\draw[->, thick] (add) -- (relu);

% skip connection
\draw[->, thick, thwsorange] (-1.8,0.3) -- (-1.8,-2.5) -- (add.west);
\draw[->, thick] (-1.8,0.8) -- (conv1);
\node[font=\scriptsize, thwsorange] at (-2.3,-1.0) {$\vx$};
\node[font=\scriptsize] at (-2.3,0.6) {input};

\draw[->, thick] (relu) -- ++(0,-0.7);
\node[font=\scriptsize] at (0.3,-3.6) {output};
\end{tikzpicture}

% \end{column}
% \begin{column}{0.04\textwidth}
% \end{column}
% \end{columns}

% \end{frame}

% % -------------------------------------------------------
% \section{Modern Details}
% % -------------------------------------------------------

% \begin{frame}{$1\times1$ Convolutions}

% \vskip 2em

% \structure{$1\times1$ conv:} applies a learned linear combination across channels at each spatial location

% \[
%     y_c[i,j] = \sum_{c'=1}^{C_{\text{in}}} w_{c,c'}\, x_{c'}[i,j] + b_c
% \]

% \vskip 1em

% \structure{Equivalent to:} FC layer applied independently at each spatial location

% \vskip 1em

% \structure{Uses:}
% \begin{itemize}
%     \item \structure{Channel reduction/expansion} --- change $C_{\text{in}}$ to $C_{\text{out}}$ cheaply (bottleneck blocks)
%     \item \structure{Non-linearity injection} --- add ReLU between channels without spatial processing
%     \item \structure{Network-in-network} --- increase model capacity without larger filters
%     \item \structure{Projection shortcuts} --- match dimensions in residual connections
% \end{itemize}

% \end{frame}

% % -------------------------------------------------------

% \begin{frame}{Depthwise Separable Convolutions}

% \vskip 2em

% \structure{Standard conv:} $C_{\text{out}} \times C_{\text{in}} \times k \times k$ parameters

% \vskip 1em

% \structure{Depthwise separable conv:} split into two steps

% \begin{columns}[t]
% \begin{column}{0.46\textwidth}
% \structure{Depthwise conv:}
% \begin{itemize}
%     \item One filter per input channel
%     \item Spatial processing only
%     \item $C_{\text{in}} \times k \times k$ parameters
% \end{itemize}
% \end{column}
% \begin{column}{0.46\textwidth}
% \structure{Pointwise conv ($1\times1$):}
% \begin{itemize}
%     \item Mix channels
%     \item $C_{\text{out}} \times C_{\text{in}}$ parameters
% \end{itemize}
% \end{column}
% \begin{column}{0.04\textwidth}
% \end{column}
% \end{columns}

% \vskip 1em

% \structure{Parameter reduction:} factor of $\approx k^2$ --- for $3\times3$: $\sim$9$\times$ fewer parameters

% \vskip 0.5em

% \structure{Used in:} MobileNet, EfficientNet --- efficient architectures for mobile/edge deployment

% \end{frame}

% % -------------------------------------------------------

% \begin{frame}{Batch Norm in CNNs}

% \vskip 2em

% \structure{BatchNorm2d:} normalises over batch \textbf{and} spatial dimensions for each channel

% \[
%     \mu_c = \frac{1}{N \cdot H \cdot W}\sum_{n,i,j} x_c^{(n)}[i,j], \qquad
%     \sigma^2_c = \frac{1}{N \cdot H \cdot W}\sum_{n,i,j}(x_c^{(n)}[i,j] - \mu_c)^2
% \]

% \begin{itemize}
%     \item One mean and variance per \textbf{channel} --- not per spatial location
%     \item Learnable $\gamma_c$ and $\beta_c$ per channel
% \end{itemize}

% \vskip 1em

% \structure{Placement:} Conv $\to$ \alert{BN} $\to$ ReLU (standard) or Conv $\to$ ReLU $\to$ BN (less common)

% \vskip 1em

% \structure{Effect in CNNs:}
% \begin{itemize}
%     \item Essential for training deep CNNs --- without it, networks beyond $\sim$10 layers are very hard to train
%     \item Allows higher learning rates
%     \item Reduces need for careful initialisation
% \end{itemize}

% \end{frame}

% % -------------------------------------------------------

% \begin{frame}{Dropout in CNNs}

% \vskip 2em

% \structure{Standard dropout} on FC layers works well in CNNs too

% \vskip 1em

% \structure{Problem with dropout on conv layers:}
% \begin{itemize}
%     \item Adjacent pixels are highly correlated
%     \item Dropping individual activations does not create enough independence
%     \item Information easily propagated through neighbouring non-dropped activations
% \end{itemize}

% \vskip 1.5em

% \structure{SpatialDropout2d:} drop entire feature maps (channels) rather than individual activations
% \begin{itemize}
%     \item Removes all spatial locations of a channel simultaneously
%     \item Forces network to use diverse feature maps
%     \item More effective regularisation for conv layers
% \end{itemize}

% \vskip 1em

% \structure{In practice:} BatchNorm often sufficient regularisation in conv layers --- dropout mainly used in FC head

% \end{frame}

% % -------------------------------------------------------
% \section{Summary}
% % -------------------------------------------------------

% \begin{frame}{Summary}

% \vskip 1em

% \begin{columns}[t]
% \begin{column}{0.48\textwidth}
% \structure{Why CNNs:}
% \begin{itemize}
%     \item Local connectivity
%     \item Parameter sharing
%     \item Translation invariance
% \end{itemize}

% \vskip 0.5em

% \structure{Building blocks:}
% \begin{itemize}
%     \item Conv: feature extraction
%     \item BatchNorm + ReLU: stability + non-linearity
%     \item Pooling / stride: downsampling
%     \item FC head: classification
% \end{itemize}

% \vskip 0.5em

% \structure{Key architectures:}
% \begin{itemize}
%     \item LeNet $\to$ AlexNet $\to$ VGG $\to$ ResNet
%     \item Residual connections: key to deep networks
% \end{itemize}
% \end{column}
% \begin{column}{0.48\textwidth}
% \structure{Modern details:}
% \begin{itemize}
%     \item $1\times1$ conv: channel mixing
%     \item Depthwise separable: efficiency
%     \item BatchNorm2d: per-channel normalisation
%     \item SpatialDropout: conv regularisation
% \end{itemize}

% \vskip 0.5em

% \structure{In PyTorch:}
% \begin{itemize}
%     \item \texttt{nn.Conv2d}
%     \item \texttt{nn.MaxPool2d}
%     \item \texttt{nn.BatchNorm2d}
%     \item \texttt{nn.AdaptiveAvgPool2d}
% \end{itemize}

% \vskip 0.5em

% \structure{Next:} L10 --- Sequence Models \& RNNs
% \end{column}
% \end{columns}

% \end{frame}

\end{document}
